% TOP_PROBLEM: {"id": 8400004, "title": "O5a — Deterministic polynomial-time integer factorization", "category_id": 3, "category": {"id": 3, "name": "graph_theory", "display_name": "Graph Theory", "description": "Problems involving graphs, networks, and their properties.", "slug": "graph-theory", "order_index": 3, "created_at": "2026-07-31T15:26:25.671Z"}, "status": "open"}
% FIRST_POSTED: null
% !TeX program = lualatex
% Compile twice: lualatex open_problems_500.tex
% All references and prose are embedded; no BibTeX database or external inputs.
\documentclass[11pt,a4paper]{article}
\usepackage[margin=24mm,headheight=14pt]{geometry}
\usepackage{fontspec}
\setmainfont{Latin Modern Roman}
\setsansfont{Latin Modern Sans}
\setmonofont{Latin Modern Mono}
\usepackage{amsmath,amssymb,mathtools}
\usepackage{unicode-math}
\setmathfont{Latin Modern Math}
\usepackage{microtype}
\usepackage{enumitem,xurl,needspace}
\usepackage[dvipsnames]{xcolor}
\usepackage{hyperref}
\hypersetup{unicode=true,colorlinks=true,linkcolor=MidnightBlue,urlcolor=MidnightBlue,
 pdftitle={Mathematical Research Notebook},pdfauthor={Source-annotated compilation},bookmarksnumbered=true}
\usepackage{bookmark}
\usepackage{tocloft}
\setlength{\cftsecnumwidth}{3em}
\cftsetpnumwidth{2.5em}
\cftsetrmarg{4em}
\usepackage{titlesec}
\titleformat{\section}{\Large\bfseries\raggedright}{\thesection}{0.7em}{}
\usepackage{fancyhdr}
\pagestyle{fancy}\fancyhf{}
\fancyhead[L]{\small\sffamily Open Problems Math | Research notebook}
\fancyhead[R]{\small 27 September 2026}
\fancyfoot[C]{\thepage}
\setlength{\parindent}{0pt}
\setlength{\parskip}{4pt plus 1pt}
\setlength{\emergencystretch}{3em}
\setcounter{tocdepth}{1}
\setcounter{secnumdepth}{1}
\setlist[itemize]{leftmargin=3em,itemsep=3pt,topsep=4pt,parsep=0pt}
\allowdisplaybreaks
\newcommand{\N}{\mathbb{N}}
\newcommand{\Z}{\mathbb{Z}}
\newcommand{\Q}{\mathbb{Q}}
\newcommand{\R}{\mathbb{R}}
\newcommand{\C}{\mathbb{C}}
\newcommand{\F}{\mathbb{F}}
\newcommand{\A}{\mathbb{A}}
\newcommand{\PP}{\mathbb P}
\newcommand{\E}{\mathbb{E}}
\newcommand{\eps}{\varepsilon}
\newcommand{\dd}{\,\mathrm d}
\newcommand{\sref}[2]{\hyperlink{source:#1:#2}{[S#2]}}
\newcommand{\eref}[2]{\hyperlink{extra:#1:#2}{[E#2]}}
% Turn the next switch off for a shorter reading copy. The quotations preserve
% every exactTarget field, including qualifications omitted from a short summary.
\newif\ifshowcatalogue
\showcataloguetrue
\newcommand{\cataloguescope}[1]{\ifshowcatalogue
 \paragraph{Exact scope in the supplied catalogue.}
 \begin{quote}\small #1\end{quote}\fi}
\usepackage{etoolbox}
\AtBeginEnvironment{itemize}{\small}
% Standalone entry; original section content is delimited for verification.
\begin{document}
\setcounter{section}{0}
%% BEGIN ORIGINAL SECTION CONTAINER
% ================================================================
% problemId: 8400004
% Canonical problem ID: 8400004
% exactTarget SHA-256: b68294af298392bbad8a4ad4a1a50595e5be9a8ee232127d2e28950b3f8cb2bf
\clearpage
\section*{\texorpdfstring{O5a — Deterministic polynomial-time integer factorization}{O5a — Deterministic polynomial-time integer factorization}}\label{8400004}
\subsection{Definitions and mathematical statement}
The binary length of $N\ge2$ is $\ell(N)=\lfloor\log_2N\rfloor+1$. A nontrivial factor is an integer $d$ with $1<d<N$ and $d\mid N$. Ask whether there exist a deterministic classical algorithm $A$ and constants $C,k$ such that
\[
 \forall N\text{ composite},\qquad A(N)=d\text{ a nontrivial factor},\qquad
 T_A(N)\le C\ell(N)^k.
\]
The promise is compositeness; the running-time bound is in the bit length, not in the value of $N$.
\par\smallskip\textit{Formulation and concept references:} \sref{8400004}{1}.
\cataloguescope{Determine whether there is a deterministic classical algorithm that, on input a composite integer N in binary, outputs a nontrivial factor of N in time polynomial in log N.}
\subsection{Short English statement}
Is there a deterministic classical method that factors every composite integer in time polynomial in its number of binary digits?
%% BEGIN RESEARCH INSERTION
\subsection{Research attempt}
\paragraph{Current frontier.}
\textit{Research check: 27 September 2026.}
Harvey's deterministic algorithm has an exponent $1/5$ in the integer's
\emph{value}, up to logarithmic factors, rather than polynomial complexity
in its bit length. \eref{8400004}{1}. Umans--Wang's November 2025 manuscript
proposes a number-theoretic hypothesis that could lower this exponent to
$1/6$; it does not claim the polynomial-time algorithm asked for here.
\sref{8400004}{2}. The research below isolates two distinct obstructions in an
order-based route: computing an exponent and deterministically finding a
base that separates prime-power components.

\paragraph{Attack route.}
Near-square search, smooth-order methods, and modular-square-root methods
all find factors on structured inputs. A global solution needs a uniform
coverage theorem for all composites. We investigate a modular-exponent
certificate because its verification and its failure mechanism can both
be written exactly. We also compare it with a geometric near-square
criterion, rather than assuming that random-base success derandomizes.

\paragraph{Attempt: factoring with a supplied exponent.}
Let $N$ be odd, with at least two distinct prime divisors, and suppose a
positive integer $M$ divisible by the Carmichael exponent $\lambda(N)$
is supplied, with bit length polynomial in $\log N$. Write $M=2^su$,
$u$ odd. For any base $a$, first compute $\gcd(a,N)$. A proper gcd is
already a factor. For a unit, compute
\[
 b_0=a^u\pmod N,\qquad b_{j+1}=b_j^2\pmod N,
 \qquad 0\leq j<s.
 \tag{17.1}
\]
Test all $\gcd(b_j-1,N)$. These computations take polynomial time in
$\log N+\log M$ by binary exponentiation and the Euclidean algorithm.

To characterize success, factor $N=\prod_{i=1}^r p_i^{e_i}$ only for the
\emph{analysis}, not as an algorithmic subroutine. Put
$t_i=v_2(p_i-1)$ and
$v_i=v_2(\operatorname{ord}_{p_i^{e_i}}(a))$. Since $u$ removes the odd
part of every component order, $b_j\equiv1\pmod{p_i^{e_i}}$ exactly when
$v_i\leq j$. If the $v_i$ are not all equal, choosing $j$ between their
minimum and maximum yields a nontrivial gcd in (17.1). Conversely if all
are equal, the chain simultaneously becomes one in every component and
none of these gcds is proper. To exclude a partial prime-power gcd,
note that each $b_j$ has two-power order: reduction from units modulo
$p_i^{e_i}$ to units modulo $p_i$ has an odd-order kernel, so a nonidentity
element of this two-primary subgroup cannot reduce to one. Thus this
algorithm succeeds on a unit
\emph{if and only if} its component two-adic order valuations differ.

For a uniformly random unit, the CRT makes the components independent.
The cyclic group of units modulo an odd prime power has two-primary
part of order $2^{t_i}$. Consequently
\[
 \Pr(v_i=0)=2^{-t_i},\qquad
 \Pr(v_i=j)=2^{j-1-t_i}\quad(1\leq j\leq t_i).
 \tag{17.2}
\]
Each atom is at most $1/2$. For $r\geq2$,
\[
 \Pr(\hbox{failure}\mid a\text{ a unit})
 =\sum_{j\geq0}\prod_{i=1}^r\Pr(v_i=j)\leq\frac12.
 \tag{17.3}
\]
Indeed, bound one of the first two factors by $1/2$, sum the other, and
bound all remaining factors by one. Sampling from $1,\ldots,N-1$ is also
safe: a nonunit has a proper gcd immediately, so overall success is at
least $1/2$. The argument establishes a randomized algorithm \emph{given}
$M$. It establishes neither a deterministic base choice nor a way to
compute $M$. The order/square-root mechanism is standard; the displayed
probability calculation is supplied here in full. \eref{8400004}{2}.

\paragraph{Attempt: the exact deterministic package still missing.}
A sufficient deterministic package would provide, for every such $N$,
in time polynomial in $\log N$, an exponent $M$ as above and a list
$\mathcal A_N$ of polynomially many bases containing either a nonunit
with proper gcd or a unit whose valuations $v_i$ differ. Equations
(17.1)--(17.3) prove sufficiency, but a density lower bound does not
construct this list: the successful set depends on the unknown factors.
Trying the first $\operatorname{poly}(\log N)$ integers is an additional
number-theoretic claim, not a consequence of (17.3).

The input exceptions do not create a loophole. Even $N$ is immediate.
Perfect powers can be recognized deterministically by testing all integer
exponents $2\leq e\leq\lfloor\log_2N\rfloor$ and computing integer
$e$-th roots by binary search; successful roots provide a factor. An odd
composite that is not a perfect power has at least two distinct prime
divisors, so the proposed package would cover all remaining inputs.
Computing the exponents and a hitting list uniformly is precisely the
unproved part of this attack.

\paragraph{Attempt: comparison with near-square search.}
For an odd semiprime $N=pq$, $p\leq q$, Fermat's method tests integers
$a\geq\lceil\sqrt N\rceil$ until $a^2-N=b^2$. The pair
\[
 a_*=(p+q)/2,\qquad b_*=(q-p)/2
\]
works. The number of candidates through this pair is at most
\[
 1+a_*-\sqrt N
 =1+\frac{(\sqrt q-\sqrt p)^2}{2}
 =1+\frac{(q-p)^2}{2(\sqrt q+\sqrt p)^2}.
 \tag{17.4}
\]
All square and square-root tests have polynomial bit cost. Thus
$q-p\leq N^{1/4}(\log N)^C$ gives polynomial time for fixed $C$, since
$(\sqrt q+\sqrt p)^2\geq4\sqrt N$. This is a proven special-case
criterion, not a generic factoring estimate. With very unequal primes,
(17.4) can be exponential in $\log N$. Restricting attention to balanced
semiprimes would silently weaken the target, and balancedness alone
still does not give the $N^{1/4}$ difference bound.

\paragraph{Stress test.}
An exponent can be short without being efficiently obtainable:
$\lambda(N)$ itself has $O(\log N)$ bits. Using its existence as though
it were computed is invalid. The easily computable integer $N!$ is a
multiple of the needed orders but has exponentially many bits in the
input length; it cannot be inserted as a free certificate. All bounds
above distinguish bit cost from arithmetic-operation count. The random
unit calculation is checked on small composites in the supplied code,
but those tests cannot certify a polynomial-size deterministic hitting
set for every input. A polynomial-time solution would not by itself imply
$\mathsf P=\mathsf{NP}$, and a negative deterministic result would not
by itself give the randomized lower bound needed in UnsolvedMath problem 30006683.

\paragraph{Outcome.}
\textbf{Outcome: Exploratory approach with unresolved gap.}
The attempt proves a complete success criterion and probability bound for
an exponent-assisted algorithm, as well as a uniform near-square special
case. It identifies exactly what a deterministic extension requires.
Neither efficient exponent production nor the necessary deterministic
coverage lemma is established. These auxiliary facts are not claimed to
be historically novel, and no general factoring complexity conclusion is
asserted.
%% END RESEARCH INSERTION
\subsection{Sources}
\begin{itemize}\raggedright
\item[\textnormal{[S1]}]\hypertarget{source:8400004:1}{} Integer factoring and modular square roots.
\url{https://www.sciencedirect.com/science/article/pii/S0022000015000768}.
{\small\textit{Catalogue formulation source.}}
\item[\textnormal{[S2]}]\hypertarget{source:8400004:2}{} A number-theoretic conjecture implying faster algorithms for polynomial factorization and integer factorization — Chris Umans and Siki Wang (2025).
\url{https://arxiv.org/abs/2511.10851}.
{\small\textit{Catalogue research/status reference; not a proof endorsement.}}
\item[\textnormal{[S3]}]\hypertarget{source:8400004:3}{} Improved Separation Between Quantum and Classical Computers for Sampling and Functional Tasks — Marshall, Aaronson and Dunjko, CCC 2025.
\url{https://drops.dagstuhl.de/storage/00lipics/lipics-vol339-ccc2025/LIPIcs.CCC.2025.5/LIPIcs.CCC.2025.5.pdf}.
{\small\textit{Catalogue research/status reference; not a proof endorsement.}}
%% BEGIN RESEARCH REFERENCES
\item[\textnormal{[E1]}]\hypertarget{extra:8400004:1}{} David Harvey, \emph{An exponent one-fifth algorithm for deterministic integer factorisation}, arXiv:2010.05450 (2020), main theorem.
\url{https://arxiv.org/abs/2010.05450}.
{\small\textit{Added research source. Deterministic value-exponent bound, not a polynomial bit-time algorithm. Consulted 27 September 2026.}}
\item[\textnormal{[E2]}]\hypertarget{extra:8400004:2}{} Peter W. Shor, \emph{Polynomial-Time Algorithms for Prime Factorization and Discrete Logarithms on a Quantum Computer}, SIAM Journal on Computing 26(5) (1997), 1484--1509; arXiv:quant-ph/9508027.
\url{https://arxiv.org/abs/quant-ph/9508027}.
{\small\textit{Added research source. Factoring through order and nontrivial square-root information; the exponent-assisted classical calculation is proved separately. Consulted 27 September 2026.}}
%% END RESEARCH REFERENCES
\end{itemize}

%% END ORIGINAL SECTION CONTAINER
\end{document}
